\documentclass[../main]{subfiles}
\begin{document}
\section{TIPOS DE LÍNEAS}
\subsection*{Líneas}
Las líneas en dibujo técnico se utilizan para diferentes propósitos como proporcionar información específica para los diseñadores, fabricantes, técnicos, arquitectos, etc. Estas líneas son un tipo de lenguaje entre los técnicos, por eso están normalizadas y debemos saberlas interpretar.

\begin{figure}[ht]
    \centering
    \includegraphics[width=\textwidth]{images/03/LINEAS.jpg}
    \caption{Tipos de líneas}
    \label{fig:my_lineas}
\end{figure}

 La combinación de los tipos de líneas con sus espesores normalizados tienen el objetivo de lograr el mejor contraste entre ellas y que la interpretación de un dibujo sea más fácil.

 Los tipos de línea a emplear en dibujo técnico vienen especificadas en la normas IRAM (Instituto de Racionalización Argentino de Materiales).De forma general, tenemos los siguientes tipos de líneas:
\begin{itemize}
 \item Línea Gruesa: Para contornos y aristas visibles.
 \item Línea Fina : Líneas de cota, líneas auxiliares de cota, líneas de ejes y líneas de rayado.
 \item Línea Fina de Trazo y Punto: Para Ejes de revolución y Simetrias.
 \item Línea fina de trazos: Contornos y Aristas Ocultas.
 \item Línea fina a Mano Alzada: Límites de vistas o cortes parcialmente interrumpidos.
 \end{itemize}
 
 La relación de anchuras entre líneas gruesas y finas no debe ser inferior a 2 (el doble o la mitad) y la anchura de la línea deberá elegirse entre la gama siguiente:

 0,25mm – 0,35mm – 0,5mm – 0,7mm – 1mm – 1,4mm y 2 mm.

 También podría ser de 0,18mm pero no suele usarse por ser demasiado fina.

\begin{figure}[ht]
    \centering
    \includegraphics[width=\textwidth]{images/03/acota1.png}
    \caption{Ejemplo de aplicación de líneas en dibujo técnico.}
    \label{fig:my_ejemplolineas}
\end{figure}

 Las líneas finas, medias y gruesas dentro de un dibujo o plano, deberán tener siempre la misma anchura, lógicamente cada una la suya propia.

\begin{tcolorbox}
Video: Tipos de línea en dibujo técnico.
\tcblower
\qrcode{https://www.youtube.com/watch?v=3VjfQQBNZPc} https://www.youtube.com/watch?v=3VjfQQBNZPc
\end{tcolorbox}

\subsection*{Actividades:}
\begin{itemize}
\item ¿Cuál es el propósito del uso de distintos tipos de línea en dibujo técnico?
\item En qué norma están definidas los tipos de línea para dibujo técnico.
\item Copiar la figura \ref{fig:my_ejemplolineas} en escala 1:1 en la siguiente página, utilizando el tipo de línea adecuado en cada caso, usar regla y respetar las medidas indicadas en milímetros.
\end{itemize}
\end{document}